\documentclass[../cours_1re.tex]{subfiles}
\begin{document}
\setcounter{chapter}{19}
\chapter{Épreuves répétées et dénombrement}
\label{ch:20}
\begin{objectif}
Modéliser jusqu'à quatre épreuves de Bernoulli indépendantes comme le demande le programme.
Le dénombrement général et la loi binomiale sont des compléments préparant la terminale.
\end{objectif}
\section{Répéter une épreuve}
Une épreuve de Bernoulli a deux issues, succès $S$ de probabilité $p$
et échec $E$ de probabilité $1-p$. Des épreuves identiques et indépendantes
conservent ces probabilités à chaque étape, quels que soient les résultats passés.
Ainsi un chemin de $n$ résultats comportant $k$ succès a pour probabilité
$p^k(1-p)^{n-k}$, par multiplication des probabilités conditionnelles.
Les cas $p=0$ ou $p=1$ se traitent directement : tout est déterministe.

Pour trois épreuves, voici tous les chemins ; aucun n'est oublié.
\begin{center}
\begin{tabular}{clc}
\toprule Nombre de succès & Chemins & Probabilité totale\\\midrule
0 & EEE & $(1-p)^3$\\
1 & SEE, ESE, EES & $3p(1-p)^2$\\
2 & SSE, SES, ESS & $3p^2(1-p)$\\
3 & SSS & $p^3$\\\bottomrule
\end{tabular}
\end{center}
Pour quatre épreuves, les nombres de chemins sont $1,4,6,4,1$.
Par exemple les chemins à deux succès correspondent aux positions
$\{1,2\},\{1,3\},\{1,4\},\{2,3\},\{2,4\},\{3,4\}$.
\begin{exemple}
Pour quatre épreuves avec $p=0{,}3$, la probabilité d'au moins un succès
est $1-0{,}7^4=0{,}7599$. Celle d'exactement deux succès est
$6\times0{,}3^2\times0{,}7^2=0{,}2646$.
Le complément évite d'additionner quatre cas pour la première question.
\end{exemple}

\section{Compter sans lister : coefficients binomiaux}
\approfondissement
Le nombre de choix de $k$ positions parmi $n$ est noté $\binom nk$.
On pose $0!=1$ et $n!=1\times2\times\cdots\times n$ pour $n\ge1$.
\begin{demonstration}
Pour choisir successivement $k$ positions distinctes et ordonnées, on a
$n(n-1)\cdots(n-k+1)$ possibilités. Chaque ensemble de $k$ positions
est compté $k!$ fois, une fois par ordre. Donc
\[
\binom nk=\frac{n!}{k!(n-k)!}\quad(0\le k\le n).
\]
Les choix de $0$ et de $n$ positions sont uniques.
Classer les choix selon qu'ils contiennent ou non la dernière position donne
$\binom nk=\binom{n-1}{k-1}+\binom{n-1}{k}$ pour $1\le k<n$.
\end{demonstration}
Dans le développement de $(a+b)^n$, choisir le terme $a$ dans $k$
facteurs donne $a^kb^{n-k}$, avec $\binom nk$ choix :
\[
(a+b)^n=\sum_{k=0}^n\binom nk a^kb^{n-k}.
\]
Cette preuve combinatoire démontre la formule du binôme.

\section{La loi binomiale comme prolongement}
Pour $n$ épreuves identiques indépendantes, le nombre $X$ de succès vérifie
\[
P(X=k)=\binom nk p^k(1-p)^{n-k},\quad 0\le k\le n,\quad 0<p<1.
\]
Les événements correspondant aux chemins sont disjoints ; leur réunion donne
$\{X=k\}$. Le binôme prouve que la somme des probabilités vaut
$(p+1-p)^n=1$. Les valeurs $p=0,1$ donnent respectivement $X=0$ et $X=n$
avec certitude. L'espérance $np$ et la variance $np(1-p)$ seront démontrées
avec les variables indicatrices aux chapitres 21 et 22.

\section{Deux précautions sur l'indépendance}
Deux tirages successifs sans remise ne sont généralement pas indépendants :
dans une urne contenant deux boules rouges et une bleue, la probabilité
du rouge passe de $2/3$ à $1/2$ après un premier rouge.
Avec remise et mélange, le modèle indépendant peut être pertinent.

\approfondissement
L'indépendance deux à deux ne suffit pas pour multiplier trois probabilités.
Sur les quatre couples équiprobables $(0,0),(0,1),(1,0),(1,1)$, posons
$A=$ « premier bit égal à $1$ », $B=$ « second bit égal à $1$ »,
$C=$ « bits différents ». Chaque événement a probabilité $1/2$ et chaque
intersection de deux a probabilité $1/4$. Pourtant
$P(A\cap B\cap C)=0\ne1/8$.
Pour une famille indépendante, la formule du produit doit valoir pour
toute sous-famille, pas seulement pour les paires.

\section{Une marche aléatoire}
À chaque succès on avance de $1$, à chaque échec on recule de $1$.
Après $n$ pas, la position est $S_n=X-(n-X)=2X-n$.
Elle a donc la même parité que $n$. Cette transformation d'une variable
connue permet de calculer sa loi, puis ses moments, sans reconstruire
un arbre complet.

\begin{finchapitre}
\section{Utilité et prolongements}
Le dénombrement transforme une explosion du nombre de chemins en une
formule utilisable. Les variables indicatrices simplifieront encore le calcul
des moyennes et des variances.

Les marches aléatoires préparent l'étude des chaînes de Markov et, après
un changement d'échelle nécessitant des théorèmes de convergence,
du mouvement brownien en probabilités avancées.
\begin{retenir}
Le modèle binomial exige des épreuves indépendantes de même probabilité.
Compter les chemins puis multiplier par leur probabilité. « Deux à deux » ne signifie pas « mutuellement ».
\end{retenir}
\end{finchapitre}
\end{document}
